\section{Open questions}
The replication settled the paper's own quantitative claims, but it raises
five forward-looking questions that neither the paper nor this reproduction
answers. These are listed in structured form in
\texttt{open\_questions.json}; summarized here:

\begin{enumerate}
    \item \textbf{Does the CEMM unified template extend to modern algorithms
          (HHL, VQE, QAOA)?} HHL is often called QPE-in-disguise; VQE/QAOA
          are variational and lack an obvious eigenphase to estimate. A
          principled inventory of which post-1998 algorithms fit the
          Hadamard/QFT $\to$ $f$-controlled-$U$ $\to$ QFT$^{-1}$ frame ---
          and which genuinely do not --- would be a natural sequel to the
          paper. \emph{Next step:} force-fit HHL, then QAOA/VQE, into the
          same numpy scaffold used here; document where the reduction breaks.

    \item \textbf{Modern query-complexity lower bounds.} The paper predates
          the polynomial method (Beals et al.) and adversary bounds
          (Ambainis, Reichardt). Its ``one query suffices'' claims for DJ/BV
          are trivially tight; Grover is known tight; Simon has an
          exponential gap. But no unified 1998-vs-modern comparison table
          exists. \emph{Next step:} produce that table (paper's claim vs.\
          today's tight bound) for all five algorithms.

    \item \textbf{QPE $4/\pi^{2}$ bound at benchmark sizes and under
          adversarial phases.} This replication verified the bound with $m=8$
          and $2000$ uniformly random $\varphi$ (min $0.4056$ vs.\ analytic
          $0.4053$). But the true worst case is a \emph{specific} $\varphi$
          (mid-bin between dyadic values), not a random draw. \emph{Next
          step:} derive the worst-case $\varphi$ analytically, evaluate at
          $m=8,12,16,20$, and confirm exact saturation of $4/\pi^{2}$.

    \item \textbf{Connection to modern quantum-supremacy circuits.} The
          paper says quantum algorithms have a common architecture. Sycamore/
          Zuchongzhi random-sampling circuits are deliberately structure-free.
          There is a real philosophical tension: is CEMM a claim about
          \emph{useful} algorithms, or about \emph{all} quantum computation?
          \emph{Next step:} operationally measure the diamond distance from a
          small random supremacy-style circuit to the nearest QPE circuit of
          equal depth.

    \item \textbf{ML-driven CEMM family classification.} A graph-neural
          network or circuit-transformer trained on labeled synthetic
          circuits from each of the five CEMM families could auto-classify
          new algorithms (HHL, VQE, QAOA, Shor-for-larger-$N$) into the
          taxonomy, potentially revealing hidden family memberships.
          \emph{Next step:} build the dataset ($\sim 10^{4}$ circuits per
          class), train the classifier, publish the confusion matrix on
          real-world algorithm implementations.
\end{enumerate}
